%% kapa: exam=1 round=27 year=2016 period=2 form=A subject=accounting subject_ko=회계학 seq=35 number=75
%% source: Q-Net 감정평가사 기출문제 (cst00309) artlSeq=5169402 file=extracted/1008/5169402/제27회 감정평가사 1차 문제/제27회 감정평가사 2교시문제(A형).pdf
%% conversion: kordoc 4.24.51 -> tools/gen1.py (auto); review=proofread-sonnet; flags=none
\begin{question}{75}
다음은 (주)감평의 20×6년도 예산자료이다. 손익분기점을 달성하기 위한 A제품의 예산판매수량은? (단, 매출배합은 변하지 않는다고 가정한다.)

{\small\centering
\begin{tabular}{|c|r|r|}\hline
구 분 & \multicolumn{1}{c|}{A제품} & \multicolumn{1}{c|}{B제품} \\ \hline
총매출액 & ￦2,100,000 & ￦2,900,000 \\ \hline
총변동원가 & 1,470,000 & 1,740,000 \\ \hline
총고정원가 & \multicolumn{2}{c|}{1,074,000} \\ \hline
판매량 & 600개 & 400개 \\ \hline
\end{tabular}\par}
\begin{choices}
  \item 240개
  \item 300개
  \item 360개
  \item 420개
  \item 480개
\end{choices}
\end{question}
\answer{3}
